% Generated draft: requires mathematical and source review.
\documentclass[11pt]{article}
\usepackage[margin=25mm]{geometry}
\usepackage{fontspec}
\setmainfont{Latin Modern Roman}
\usepackage{amsmath,amssymb,mathtools,unicode-math}
\setmathfont{Latin Modern Math}
\usepackage{enumitem,needspace,xurl,hyperref}
\setlength{\parindent}{0pt}
\setlength{\parskip}{5pt}
\setlength{\emergencystretch}{3em}
\newcommand{\N}{\mathbb{N}}
\newcommand{\Z}{\mathbb{Z}}
\newcommand{\Q}{\mathbb{Q}}
\newcommand{\R}{\mathbb{R}}
\newcommand{\C}{\mathbb{C}}
\newcommand{\F}{\mathbb{F}}
\newcommand{\A}{\mathbb{A}}
\newcommand{\PP}{\mathbb{P}}
\newcommand{\E}{\mathbb{E}}
\newcommand{\eps}{\varepsilon}
\newcommand{\dd}{\,\mathrm{d}}
\begin{document}
\section*{UnsolvedMath ID: 701420}
\textbf{Source problem code:} catalogue entry
\subsection{Definitions and mathematical statement}
For a finite group \(H\), its rank \(\operatorname{rk}(H)\) is the least nonnegative integer \(s\) such that every subgroup of \(H\) can be generated by at most \(s\) elements. A group is soluble if its derived series terminates at the trivial subgroup. Let \(F(H)\) denote the Fitting subgroup of a finite group \(H\), namely its largest nilpotent normal subgroup. Define the upper Fitting series by
\[
F_0(H)=1,\qquad F_{i+1}(H)/F_i(H)=F\bigl(H/F_i(H)\bigr)\quad(i\geq 0).
\]
Let \(G\) be a finite soluble group and let \(A\leq\operatorname{Aut}(G)\) be a finite soluble group with \(\gcd(|A|,|G|)=1\). Set
\[
C_G(A)=\{g\in G:\alpha(g)=g\text{ for every }\alpha\in A\},\qquad r=\operatorname{rk}(C_G(A)).
\]
Write \(|A|=p_1\cdots p_l\), where the primes are counted with multiplicity; for \(|A|=1\), use the empty product and \(l=0\).

The question is whether there is a universal function \(f:\mathbb Z_{\geq0}\to\mathbb Z_{\geq0}\) growing at most linearly, together with a function
\[
B:\mathbb Z_{\geq1}\times\mathbb Z_{\geq0}\longrightarrow\mathbb Z_{\geq0},
\]
such that, for every such pair \((G,A)\),
\[
\operatorname{rk}\bigl(G/F_{f(l)}(G)\bigr)\leq B(|A|,r).
\]
Here growing at most linearly means that there are universal nonnegative constants \(c,d\) such that \(f(l)\leq cl+d\) for every \(l\geq0\). Thus \(f\) is independent of \(G\), \(A\), and \(r\), and the rank bound depends only on \(|A|\) and \(r\). The source says ``linear function'' without specifying whether an additive constant is allowed; the formulation above records the at-most-linear growth interpretation.

The supplied source also reports that an exponential function \(f\) with the required rank-bound property was obtained in the cited Khukhro--Mazurov reference. It further attributes to Hartley--Isaacs the assertion that there is a function \(D\), depending only on its displayed arguments, such that
\[
\bigl|G/F_{2l+1}(G)\bigr|\leq D\bigl(|C_G(A)|,|A|\bigr)
\]
under the same hypotheses. These are supplied background assertions, not independently verified results in this draft.
\subsection{Short English statement}
Does a universal number of Fitting-series steps growing linearly with the number of prime factors of \(|A|\), counted with multiplicity, leave a quotient whose rank is bounded only by \(|A|\) and the rank of the common fixed-point subgroup? The hypotheses require both groups to be finite and soluble and their orders to be coprime. The registry labels the problem open; this draft does not independently verify that status.
\subsection{Sources}
\begin{itemize}
\item Local UnsolvedMath registry: permanent ID \(701420\), exact source code \texttt{catalogue entry}, title ``Kourovka Notebook Problem 17.119.'' The supplied record labels the problem open and attributes it to E. I. Khukhro, proposed in 2010.
\item \emph{The Kourovka Notebook (Unsolved Problems in Group Theory)}, original problem 17.119, 17th issue (2010); the supplied background also identifies the 21st-issue edition (2026), \texttt{arXiv:1401.0300v48}. The supplied collection description names E. I. Khukhro and V. D. Mazurov as editors.
\item Supplied reference for the exponential-index rank bound: E. Khukhro and V. Mazurov, \emph{Groups St. Andrews 2005}, vol. II, Cambridge Univ. Press, 2007, pp. 564--585. The record does not give an article title.
\item Supplied attribution for the bound on \(|G/F_{2l+1}(G)|\): Hartley--Isaacs. No bibliographic details for this result are supplied.
\end{itemize}
\end{document}
